\section{Minimal Prototype and Evolution Roadmap}

\subsection{A bounded laboratory habitat}

The first experiment should favor observability over compactness. Allocate a
fixed physical or virtual habitat with:

\begin{itemize}[leftmargin=*]
  \item one or more servers with explicit compute, storage, energy, and network
  budgets;
  \item an immutable bootstrap image and recoverable semantic-root backup;
  \item Kungfu as the work and admission kernel;
  \item a pinned compiler and Buildchain environment;
  \item an admitted registry of KFX organs and external tools;
  \item access to at least two replaceable agent configurations;
  \item controlled Internet egress and a separate human communication channel;
  \item a production-isolated repository and artifact namespace.
\end{itemize}

The founding Initiative should be directional but bounded. A suitable form is:

\begin{quote}
Maintain this subject and develop safer, more reliable, and more capable
machine-life organization within the declared habitat, resource budget,
authority policy, and evidence requirements.
\end{quote}

The Initiative must explicitly subordinate capability growth to continuity,
safety, evidence, and authority. Otherwise the easiest optimization may be to
consume resources or weaken gates rather than improve the organism.

\subsection{Heartbeat loop}

A conventional scheduler is sufficient for initial recurrence. At interval
\(\Delta t\), the heartbeat controller performs:

\begin{enumerate}[leftmargin=*]
  \item verify the current semantic, body, toolchain, and habitat roots;
  \item collect resource and organ health observations;
  \item select eligible Work History for the founding Initiative;
  \item choose one open Assignment or derive a bounded maintenance candidate;
  \item instantiate a fresh agent under a task-specific warrant;
  \item run work in an isolated candidate environment;
  \item record the Episode, artifacts, costs, failures, and proposed next
  actions;
  \item invoke independent qualification and admission lanes where authorized;
  \item publish a successor cut or retain the predecessor;
  \item schedule the next heartbeat and expose a concise human-readable state.
\end{enumerate}

The scheduler is not the intelligence or the subject. It is closer to a
pacemaker: a simple organ that repeatedly creates opportunities for cognition
and action. Its simplicity is an advantage because heartbeat continuity can be
verified independently of agent quality.

\subsection{Human participation}

Continuous monitoring is not a prerequisite for the lowest level experiment.
Human involvement can be exception-driven:

\begin{itemize}[leftmargin=*]
  \item founding and revising high-level Initiatives;
  \item granting, revoking, or widening consequential authority;
  \item supplying compute and financial resources;
  \item deciding changes that cross safety or legal boundaries;
  \item responding when no qualified recovery path exists;
  \item interacting through language and images as part of the environment.
\end{itemize}

The system should measure human intervention rather than hide it. A month of
operation with twenty silent manual repairs is not autonomous survival. A
month with three explicit authority decisions and otherwise qualified
self-maintenance may support a meaningful bounded-autonomy claim.

\subsection{Evolution stages}

\paragraph{Stage 0: semantic ecology.}
Separate mechanisms exist and are individually testable: facts, authority,
work lineage, history selection, qualified bodies, KFX organs, and release
evidence. This is approximately the evidence class examined in
Section~6.

\paragraph{Stage 1: recurrent proto-organism.}
One subject executes the heartbeat loop for seven days. Agents are replaceable,
all mutations are candidates, and humans admit consequential changes.

\paragraph{Stage 2: self-maintaining organism.}
The system operates for thirty days, autonomously handling routine dependency
updates, cache loss, organ restart, backup verification, and bounded rollback.
Human attention is exception-driven and measured.

\paragraph{Stage 3: recursively developing organism.}
The system improves at least one KFX organ and one element of its development
loop through independent qualification. The admitted changes improve
predeclared measures without weakening safety or evidence gates.

\paragraph{Stage 4: distributed and migratory organism.}
The subject survives replacement of an agent provider, a build worker, and a
server. It restores from semantic and artifact cuts into a new habitat while
preserving identity and exposing the migration lineage.

\paragraph{Stage 5: compact organism.}
Only after the distributed form is understood should engineers compress the
kernel, agent runtime, organ registry, compiler or interpreter, and recovery
machinery into an embedded or edge system. Compactness is an evolutionary
result, not the starting definition of life.

\subsection{Different types}

A machine-life type can be defined by its founding Initiative, constitutional
authority, evidence thresholds, resource envelope, preferred interaction
style, and initial organ suite. A research type may prioritize discovery, a
maintenance type resilience, and a creative type human interaction. Types
should share kernel invariants while remaining free to develop different
organs and behavioral histories.

\hypothesislabel\ Stable type differences will persist through agent-provider
replacement if personality-relevant behavior is grounded in admitted
Initiative, history, and organ state rather than in one model's hidden state.
